\chapter{Evolution Equations}
\label{chap:cz-evolution-equations}

Throughout, $M$ is an $n$-dimensional Riemannian manifold, repeated
indices are summed, and
\[
R^k_{ijl}=\partial_i\Gamma^k_{jl}-\partial_j\Gamma^k_{il}
+\Gamma^k_{ip}\Gamma^p_{jl}-\Gamma^k_{jp}\Gamma^p_{il},\qquad
R_{ijkl}=g_{kp}R^p_{ijl}.
\]
Thus $R_{ijkl}=-R_{jikl}=-R_{ijlk}=R_{klij}$ and
\[
R_{ijkl}+R_{jkil}+R_{kijl}=0,\qquad
\nabla_mR_{ijkl}+\nabla_iR_{jmkl}+\nabla_jR_{mikl}=0.
\]
Write $R_{ik}=g^{jl}R_{ijkl}$, $R=g^{ij}R_{ij}$, and
$\Delta=g^{pq}\nabla_p\nabla_q$.  The commutator convention is
\[
[\nabla_i,\nabla_j]v^l=R^l_{ijk}v^k,\qquad
[\nabla_i,\nabla_j]v_k=R_{ijkl}g^{lm}v_m.
\]

\section{The Ricci Flow}

The Ricci flow, introduced by Hamilton \cite{Ha82}, is
\[
\frac{\partial g_{ij}}{\partial t}=-2R_{ij}.
\]
An Einstein metric $R_{ij}=\lambda g_{ij}$ evolves by
$g(t)=(1-2\lambda t)g(0)$.  Hence it is stationary when
$\lambda=0$, shrinking for $\lambda>0$, and expanding for
$\lambda<0$.

A steady Ricci soliton satisfies
\[
2R_{ij}+g_{ik}\nabla_jV^k+g_{jk}\nabla_iV^k=0,
\]
and a gradient soliton satisfies
\[
R_{ij}+\nabla_i\nabla_jf-\lambda g_{ij}=0.
\]
The cases $\lambda>0$, $\lambda=0$, and $\lambda<0$ are shrinking,
steady, and expanding, respectively.  The case of constant $f$ is
Einstein.

\begin{proposition}[compact steady and expanding gradient solitons are Einstein]
  \label{prop:cz-compact-steady-and-expanding-gradient-solitons-are-einstein}
  \source{cao-zhu}
  \dcref{Proposition 1.1.1}
  \group{Cao--Zhu The Ricci Flow}
On a compact $n$-manifold, every gradient steady or expanding Ricci
soliton is Einstein.
\end{proposition}
\begin{proof}
\sketch For a steady soliton \cite{Ha95F, CK04}, tracing and differentiating the soliton
equation give $R+\Delta f=0$ and $R+|\nabla f|^2=C$.  Compactness
gives $C=0$, hence $\int_M|\nabla f|^2\dd V=0$.  For an expanding
soliton, the corresponding identities are
$R+\Delta f=n\lambda$ and $R+|\nabla f|^2-2\lambda f=C$ with
$\lambda<0$; evaluating their difference at the extrema of $f$
forces $f$ to be constant.
\end{proof}

\section{Short-time Existence and Uniqueness}

For a fixed background metric $\mathring g$, put
\[
V_i=g_{ik}g^{jl}(\Gamma^k_{jl}-\mathring\Gamma^k_{jl}).
\]
The DeTurck-modified flow \cite{De} is
\[
\frac{\partial g_{ij}}{\partial t}=-2R_{ij}+\nabla_iV_j+\nabla_jV_i,
\qquad g(0)=\mathring g.
\]

\begin{lemma}[strict parabolicity of the DeTurck-modified flow]
  \label{lem:cz-strict-parabolicity-of-the-deturck-modified-flow}
  \source{cao-zhu}
  \dcref{Lemma 1.2.1}
  \group{Cao--Zhu Short-time Existence and Uniqueness}
The DeTurck-modified flow is strictly parabolic.
\end{lemma}
\begin{proof}
\sketch Its principal part is $g^{kl}\partial_k\partial_lg_{ij}$,
so its symbol at a nonzero covector $\zeta$ is
$(g^{kl}\zeta_k\zeta_l)\widetilde g_{ij}$.
\end{proof}

\begin{theorem}[short-time existence and uniqueness on compact manifolds]
  \label{thm:cz-short-time-existence-and-uniqueness-on-compact-manifolds}
  \source{cao-zhu}
  \dcref{Theorem 1.2.2}
  \group{Cao--Zhu Short-time Existence and Uniqueness}
For every compact Riemannian manifold $(M,g_0)$, there is $T>0$ and
a unique smooth solution of $\partial_tg=-2\Ric(g)$ on
$M\times[0,T)$ with $g(0)=g_0$.
\end{theorem}
\begin{proof}
\sketch Solve the strictly parabolic DeTurck flow \cite{Ha82, De}.  The associated
harmonic-map heat flow supplies diffeomorphisms recovering Ricci flow.
Applying the same construction to two Ricci flows reduces uniqueness to
parabolic uniqueness and then to uniqueness of the diffeomorphism ODE.
\end{proof}

\begin{theorem}[short-time existence on complete noncompact manifolds]
  \label{thm:cz-short-time-existence-on-complete-noncompact-manifolds}
  \source{cao-zhu}
  \dcref{Theorem 1.2.3}
  \group{Cao--Zhu Short-time Existence and Uniqueness}
If $(M,g_0)$ is complete, noncompact, $n$-dimensional, and has bounded
curvature, then for some $T>0$ it admits a smooth Ricci flow on
$M\times[0,T]$ with $g(0)=g_0$ and uniformly bounded curvature
\cite{Sh89}.
\end{theorem}

\begin{theorem}[bounded-curvature uniqueness on complete noncompact manifolds]
  \label{thm:cz-bounded-curvature-uniqueness-on-complete-noncompact-manifolds}
  \source{cao-zhu}
  \dcref{Theorem 1.2.4}
  \group{Cao--Zhu Short-time Existence and Uniqueness}
Two bounded-curvature Ricci flows on $M\times[0,T]$ with the same
complete noncompact bounded-curvature initial metric agree identically
\cite{CZ05U}.
\end{theorem}

\section{Evolution of Curvatures}

\begin{proposition}[evolution of the Riemann curvature tensor]
  \label{prop:cz-evolution-of-the-riemann-curvature-tensor}
  \source{cao-zhu}
  \dcref{Proposition 1.3.1}
  \group{Cao--Zhu Evolution of Curvatures}
Under Ricci flow,
\begin{align*}
\partial_tR_{ijkl}={}&\Delta R_{ijkl}
+2(B_{ijkl}-B_{ijlk}-B_{iljk}+B_{ikjl})\\
&-g^{pq}(R_{pjkl}R_{qi}+R_{ipkl}R_{qj}+R_{ijpl}R_{qk}+R_{ijkp}R_{ql}),
\end{align*}
where $B_{ijkl}=g^{pr}g^{qs}R_{piqj}R_{rksl}$.
\cite{Ha82}
\end{proposition}
\begin{proof}
\sketch Differentiate the connection and curvature in normal
coordinates.  Use the second Bianchi identity and commutation of
covariant derivatives to replace second derivatives of $\Ric$ by
$\Delta\Rm$ and the stated quadratic terms.
\end{proof}

\begin{corollary}[evolution of the Ricci tensor]
  \label{cor:cz-evolution-of-the-ricci-tensor}
  \source{cao-zhu}
  \dcref{Corollary 1.3.2}
  \group{Cao--Zhu Evolution of Curvatures}
\[
\partial_tR_{ik}=\Delta R_{ik}+2g^{pr}g^{qs}R_{piqk}R_{rs}
-2g^{pq}R_{pi}R_{qk}.
\]
\end{corollary}
\begin{proof}
\sketch Contract Proposition~1.3.1 and use the Bianchi identity to
cancel the contracted $B$-terms.
\end{proof}

\begin{corollary}[evolution of scalar curvature]
  \label{cor:cz-evolution-of-scalar-curvature}
  \source{cao-zhu}
  \dcref{Corollary 1.3.3}
  \group{Cao--Zhu Evolution of Curvatures}
\[
\partial_tR=\Delta R+2|\Ric|^2.
\]
\end{corollary}
\begin{proof}
\sketch Trace the Ricci evolution equation and use
$\partial_tg^{ij}=2R^{ij}$.
\end{proof}

Choose an evolving orthonormal frame $F_a=F_a^i\partial_i$ with
$\partial_tF_a^i=g^{ij}R_{jk}F_a^k$.  In this frame the curvature
operator is the symmetric form on $\Lambda^2$ determined by
\[
R_{abcd}=M_{\alpha\beta}\varphi_{ab}^\alpha\varphi_{cd}^\beta,
\qquad
[\varphi^\alpha,\varphi^\beta]=C^{\alpha\beta}_\gamma\varphi^\gamma.
\]

\begin{proposition}[evolution of the curvature operator]
  \label{prop:cz-evolution-of-the-curvature-operator}
  \source{cao-zhu}
  \dcref{Proposition 1.3.4}
  \group{Cao--Zhu Evolution of Curvatures}
Under Ricci flow,
\[
\partial_tM_{\alpha\beta}=\Delta M_{\alpha\beta}
+(M^2)_{\alpha\beta}+M^\#_{\alpha\beta},
\]
where $(M^2)_{\alpha\beta}=M_{\alpha\gamma}M_{\beta\gamma}$ and
\[
M^\#_{\alpha\beta}=C_\alpha^{\gamma\eta}C_\beta^{\delta\theta}
M_{\gamma\delta}M_{\eta\theta}.
\]
\cite{Ha86}
\end{proposition}
\begin{proof}
\sketch Express the quadratic part of the Riemann-tensor evolution in
an orthonormal basis of $\Lambda^2$.  The two Bianchi reductions yield
the operator square and the Lie-algebra square, respectively.
\end{proof}

In dimension $3$, $M^\#$ is the adjugate matrix
$\det(M){}^tM^{-1}$.  In dimension $4$, for
$\Lambda^2=\Lambda^2_+\oplus\Lambda^2_-$ and
$M=\left(\begin{smallmatrix}A&B\\{}^tB&C\end{smallmatrix}\right)$,
\[
M^\#=2\left(\begin{smallmatrix}A^\#&B^\#\\{}^tB^\#&C^\#\end{smallmatrix}\right),
\qquad \tr A=\tr C=\frac12R.
\]
If the Ricci tensor is diagonal in the chosen frame, then $B$ is
diagonal; in particular, $B=0$ for an Einstein four-manifold.

\section{Derivative Estimates}

\begin{theorem}[global curvature derivative estimates]
  \label{thm:cz-global-curvature-derivative-estimates}
  \source{cao-zhu}
  \dcref{Theorem 1.4.1}
  \group{Cao--Zhu Derivative Estimates}
There are constants $C_m$, depending only on $m$ and the dimension,
such that a complete Ricci flow satisfying $|\Rm|\le M$ up to
$0<t\le M^{-1}$ satisfies
\[
|\nabla\Rm|\le C_1M t^{-1/2},\qquad
|\nabla^m\Rm|\le C_mM t^{-m/2}.
\]
All norms are taken in the evolving metric \cite{Sh89}.
\end{theorem}
\begin{proof}
  \uses{prop:cz-evolution-of-the-riemann-curvature-tensor}
\sketch The equations $\partial_t\Rm=\Delta\Rm+\Rm*\Rm$ and its
covariant derivatives give differential inequalities for
$|\nabla^m\Rm|^2$.  Apply the maximum principle first to
$t|\nabla\Rm|^2+A|\Rm|^2$, then inductively to weighted sums of
successive derivative norms.
\end{proof}

\begin{theorem}[local curvature derivative estimates]
  \label{thm:cz-local-curvature-derivative-estimates}
  \source{cao-zhu}
  \dcref{Theorem 1.4.2}
  \group{Cao--Zhu Derivative Estimates}
There are positive constants $\theta,C_k$, depending only on the
dimension, with the following property.  If $|\Rm|\le M$ on
$U\times[0,\theta/M]$, $\overline B_0(p,r)\subset U$, and
$t\le\theta/M$, then
\[
|\nabla\Rm(p,t)|^2\le C_1M^2\left(r^{-2}+t^{-1}+M\right),
\]
and
\[
|\nabla^k\Rm(p,t)|^2\le C_kM^2\left(r^{-2k}+t^{-k}+M^k\right).
\]
\end{theorem}
\begin{proof}
\sketch Use a cutoff $\varphi$ supported in $B_0(p,r)$ and compare
a Bernstein quantity for $|\nabla\Rm|^2$ with a barrier.  The two
following propositions control the barrier over the required time.
The maximum principle gives the first estimate, and the same argument
applied to adjacent derivative orders proves the general case \cite{Sh89, Ha95F}.
\end{proof}

\begin{proposition}[cutoff barrier is a strict supersolution]
  \label{prop:cz-cutoff-barrier-is-a-strict-supersolution}
  \source{cao-zhu}
  \dcref{Claim 1}
  \group{Cao--Zhu Derivative Estimates}
Let
\[
H=\frac{(12+4\sqrt n)A^2}{\varphi^2}+\frac1t+M.
\]
If $|\nabla\varphi|^2\le2A^2$ and
$\varphi|\nabla^2\varphi|\le2A^2$, then
\[
\partial_tH>\Delta H-H^2+M^2.
\]
\end{proposition}
\begin{proof}
\sketch Differentiate $\varphi^{-2}$ and insert the two assumed
cutoff bounds; the $\varphi^{-4}$ and $t^{-2}$ terms in $H^2$ give
the strict inequality.
\end{proof}

\begin{proposition}[short-time control of cutoff derivatives]
  \label{prop:cz-short-time-control-of-cutoff-derivatives}
  \source{cao-zhu}
  \dcref{Claim 2}
  \group{Cao--Zhu Derivative Estimates}
For $\theta>0$ sufficiently small in terms of $b,B,A$, if
$r\le\theta/\sqrt M$, $t\le\theta/M$, and $F\le H$, then
\[
|\nabla\varphi|^2\le2A^2,qquad
\varphi|\nabla^2\varphi|\le2A^2.
\]
\end{proposition}
\begin{proof}
\sketch Evolve the first and second covariant derivatives of the
time-independent cutoff in the evolving orthonormal frame.  The bound
$F\le H$ controls $|\nabla\Rm|$, and Gronwall estimates give both
inequalities for sufficiently small $\theta$.
\end{proof}

\section{Variational Structure and Dynamic Property}

Assume in this section that $M$ is compact.

\begin{definition}[Ricci breathers]
  \label{def:cz-ricci-breathers}
  \source{cao-zhu}
  \dcref{Definition 1.5.1}
  \group{Cao--Zhu Variational Structure and Dynamic Property}
A Ricci flow $g(t)$ is a breather if, for some $t_1<t_2$ and
$\alpha>0$, the metrics $\alpha g(t_1)$ and $g(t_2)$ differ by a
diffeomorphism.  It is steady, shrinking, or expanding according as
$\alpha=1$, $\alpha<1$, or $\alpha>1$.
\end{definition}

Define Perelman's functional \cite{P1}
\[
\mathcal F(g,f)=\int_M(R+|\nabla f|^2)e^{-f}\dd V.
\]
\begin{lemma}[first variation of Perelman's F-functional]
  \label{lem:cz-first-variation-of-perelmans-f-functional}
  \source{cao-zhu}
  \dcref{Lemma 1.5.2}
  \group{Cao--Zhu Variational Structure and Dynamic Property}
For $v_{ij}=\delta g_{ij}$, $h=\delta f$, and $v=g^{ij}v_{ij}$,
\[
\delta\mathcal F(v,h)=\int_M\left[-v_{ij}(R_{ij}+\nabla_i\nabla_jf)
+\left(\frac v2-h\right)(2\Delta f-|\nabla f|^2+R)\right]e^{-f}\dd V.
\]
\end{lemma}
\begin{proof}
\sketch Use
$\delta R=-\Delta v+\nabla_i\nabla_jv_{ij}-R_{ij}v_{ij}$,
differentiate the weighted volume form, and integrate by parts against
$e^{-f}$.
\end{proof}

\begin{proposition}[monotonicity of Perelman's F-functional]
  \label{prop:cz-monotonicity-of-perelmans-f-functional}
  \source{cao-zhu}
  \dcref{Proposition 1.5.3}
  \group{Cao--Zhu Variational Structure and Dynamic Property}
If
\[
\partial_tg_{ij}=-2R_{ij},\qquad
\partial_tf=-\Delta f+|\nabla f|^2-R,
\]
then
\[
\frac{\dd}{\dd t}\mathcal F(g(t),f(t))
=2\int_M|R_{ij}+\nabla_i\nabla_jf|^2e^{-f}\dd V,
\]
and $\int_Me^{-f}\dd V$ is constant.  Thus $\mathcal F$ is
nondecreasing, strictly unless the flow is a steady gradient soliton.
\end{proposition}
\begin{proof}
  \uses{lem:cz-first-variation-of-perelmans-f-functional}
\sketch Substitute the coupled flow in the first-variation formula
and integrate by parts using the contracted Bianchi identity.  The
remaining terms form the indicated square.  Moreover,
$\partial_t(e^{-f}\dd V)=-\Delta(e^{-f})\dd V$.
\end{proof}

Set
\[
\lambda(g)=\inf\left\{\mathcal F(g,f):\int_Me^{-f}\dd V=1\right\},
\qquad
\bar\lambda(g)=\lambda(g)\Vol(g)^{2/n}.
\]
Equivalently, $\lambda(g)$ is the first eigenvalue of $-4\Delta+R$.

\begin{corollary}[lambda monotonicity and steady breathers]
  \label{cor:cz-lambda-monotonicity-and-steady-breathers}
  \source{cao-zhu}
  \dcref{Corollary 1.5.4}
  \group{Cao--Zhu Variational Structure and Dynamic Property}
Along Ricci flow, $\lambda(g(t))$ is nondecreasing and is strictly
increasing unless the flow is a steady gradient soliton.  Every steady
breather is a steady gradient soliton.
\end{corollary}
\begin{proof}
\sketch Evolve a minimizing terminal value of $f$ backwards by the
conjugate equation and apply $\mathcal F$ monotonicity.  Diffeomorphism
invariance of $\lambda$ forces equality on a steady breather.
\end{proof}

\begin{corollary}[scale-invariant lambda monotonicity and expanding breathers]
  \label{cor:cz-scale-invariant-lambda-monotonicity-and-expanding-breathers}
  \source{cao-zhu}
  \dcref{Corollary 1.5.5}
  \group{Cao--Zhu Variational Structure and Dynamic Property}
Where $\bar\lambda(g(t))\le0$, it is nondecreasing and is strictly
increasing unless the flow is a gradient expanding soliton.  Every
expanding breather is a gradient expanding soliton.
\end{corollary}
\begin{proof}
\sketch Differentiate $\mathcal F\Vol^{2/n}$ along a backward
minimizer.  The derivative is bounded below by the sum of a traceless
soliton square and a weighted variance, hence is nonnegative.  On an
expanding breather, scale and diffeomorphism invariance force equality.
\end{proof}

\begin{proposition}[steady and expanding breathers are Einstein]
  \label{prop:cz-steady-and-expanding-breathers-are-einstein}
  \source{cao-zhu}
  \dcref{Proposition 1.5.6}
  \group{Cao--Zhu Variational Structure and Dynamic Property}
On a compact manifold, every steady or expanding breather is Einstein.
\end{proposition}
\begin{proof}
\sketch The preceding corollaries make the breather a gradient steady
or expanding soliton; Proposition~1.1.1 then applies.
\end{proof}

For $\tau>0$, define Perelman's functional \cite{P1}
\[
\mathcal W(g,f,\tau)=\int_M[\tau(R+|\nabla f|^2)+f-n]
(4\pi\tau)^{-n/2}e^{-f}\dd V.
\]
It is invariant under diffeomorphisms and under
$(g,f,\tau)\mapsto(a\varphi^*g,\varphi^*f,a\tau)$.

\begin{lemma}[first variation of Perelman's W-functional]
  \label{lem:cz-first-variation-of-perelmans-w-functional}
  \source{cao-zhu}
  \dcref{Lemma 1.5.7}
  \group{Cao--Zhu Variational Structure and Dynamic Property}
For $v_{ij}=\delta g_{ij}$, $h=\delta f$, $\eta=\delta\tau$, and
$v=g^{ij}v_{ij}$,
\begin{align*}
\delta\mathcal W={}&-\int_M\tau v_{ij}
\left(R_{ij}+\nabla_i\nabla_jf-\frac{1}{2\tau}g_{ij}\right)
(4\pi\tau)^{-n/2}e^{-f}\dd V\\
&+\int_M\left(\frac v2-h-\frac{n\eta}{2\tau}\right)
[\tau(R+2\Delta f-|\nabla f|^2)+f-n-1]
(4\pi\tau)^{-n/2}e^{-f}\dd V\\
&+\int_M\eta\left(R+|\nabla f|^2-\frac{n}{2\tau}\right)
(4\pi\tau)^{-n/2}e^{-f}\dd V.
\end{align*}
\end{lemma}
\begin{proof}
  \uses{lem:cz-first-variation-of-perelmans-f-functional}
\sketch Differentiate $\mathcal W$, use the scalar-curvature
variation from the $\mathcal F$ calculation, and integrate the weighted
derivative terms by parts.
\end{proof}

\begin{proposition}[monotonicity of Perelman's W-functional]
  \label{prop:cz-monotonicity-of-perelmans-w-functional}
  \source{cao-zhu}
  \dcref{Proposition 1.5.8}
  \group{Cao--Zhu Variational Structure and Dynamic Property}
If
\[
\partial_tg_{ij}=-2R_{ij},\qquad
\partial_tf=-\Delta f+|\nabla f|^2-R+\frac{n}{2\tau},\qquad
\partial_t\tau=-1,
\]
then
\[
\frac{\dd}{\dd t}\mathcal W
=2\tau\int_M\left|R_{ij}+\nabla_i\nabla_jf-\frac{1}{2\tau}g_{ij}\right|^2
(4\pi\tau)^{-n/2}e^{-f}\dd V.
\]
Also $\int_M(4\pi\tau)^{-n/2}e^{-f}\dd V$ is constant.  Therefore
$\mathcal W$ is nondecreasing, strictly unless the flow is a
shrinking gradient soliton.
\end{proposition}
\begin{proof}
  \uses{lem:cz-first-variation-of-perelmans-w-functional}
\sketch Insert the coupled flow in the first-variation identity and
integrate by parts to complete the soliton square.  The normalized
measure evolves by a Laplacian, so its integral is constant.
\end{proof}

Define
\[
\mu(g,\tau)=\inf\left\{\mathcal W(g,f,\tau):
(4\pi\tau)^{-n/2}\int_Me^{-f}\dd V=1\right\},
\]
and let $\nu(g)$ be the infimum of $\mathcal W(g,f,\tau)$ over
$f$, $\tau>0$, and the same normalization.  A minimizer for $\mu$ is
smooth and satisfies \cite{Ro}
\[
\tau(2\Delta f-|\nabla f|^2+R)+f-n=\mu(g,\tau).
\]

\begin{corollary}[mu monotonicity and shrinking breathers]
  \label{cor:cz-mu-monotonicity-and-shrinking-breathers}
  \source{cao-zhu}
  \dcref{Corollary 1.5.9}
  \group{Cao--Zhu Variational Structure and Dynamic Property}
$\mu(g(t),\tau-t)$ is nondecreasing along Ricci flow and is strictly
increasing unless the flow is a shrinking gradient soliton.  Every
shrinking breather is a shrinking gradient soliton.
\end{corollary}
\begin{proof}
  \uses{prop:cz-monotonicity-of-perelmans-w-functional}
\sketch Evolve a terminal minimizer backwards by the conjugate heat
equation and apply $\mathcal W$ monotonicity.  For a shrinking breather,
choose $\tau$ so that $\alpha(\tau-t_1)=\tau-t_2$; scaling and
diffeomorphism invariance turn the monotonicity inequality into equality.
\end{proof}

\begin{remark}[examples of nontrivial Ricci solitons]
  \label{rem:cz-examples-of-nontrivial-ricci-solitons}
  \dcref{Remark 1.1.2}
  \group{Cao--Zhu The Ricci Flow}
  \source{cao-zhu}
By contrast, there do exist nontrivial compact gradient shrinking
Ricci solitons (see Koiso \cite{Ko}, Cao \cite{Cao94} and Wang-Zhu
\cite{WZ} ). Also, there exist complete noncompact steady gradient
Ricci solitons that are not Ricci flat. In two dimensions Hamilton
\cite{Ha88} wrote down the first such example on $\mathbb R^2$,
called the {\bf cigar soliton}\index{soliton!cigar}, where the
metric is given by \be
ds^2=\frac {dx^2 +dy^2} {1+x^2+y^2}, %\eqno(1.1.18)
\ee and the vector field is radial, given by $V=-\partial/\partial
r=-(x\partial/\partial x+y\partial/\partial y).$ This metric has
positive curvature and is asymptotic to a cylinder of finite
circumference $2\pi$ at $\infty$. Higher dimensional examples were
found by Robert Bryant \cite{BR} on $\mathbb R^n$ in the
Riemannian case, and by the first author \cite{Cao94} on $\mathbb
C^n$ in the K\"ahler case. These examples are complete,
rotationally symmetric, of positive curvature and found by solving
certain nonlinear ODE (system). Noncompact expanding solitons were
also constructed by the first author \cite{Cao94}. More recently,
Feldman, Ilmanen and Knopf \cite{FIM} constructed new examples of
noncompact shrinking and expanding K\"ahler-Ricci solitons.
\end{remark}
